\section{Theoretical Analysis and Calculations} \label{apdx:theory}
In the following lemma, we prove a point-wise lower bound on the expected improvement before stopping for our recommended pairing of PBGI/LogEIPC stopping rule paired with either PBGI or LogEIPC acquisition function.\footnote{In fact, any acquisition function (e.g., expected improvement-cost (EIC) \citep{hu2025adjusted}) that ensures the one-step expected improvement is worth the evaluation cost before the stopping time $\tau$ can apply here.}

\LemEIlb*

\begin{proof}
While stopping has not occurred, meaning $t<\tau$, by the stopping criteria definition we have $\max_{x\in X}\alpha_{t}^{\mathrm{EI}}(x)/c(x)\ge1$.
Hence, there exists at least one point with
$\alpha_{t}^{\mathrm{EI}}(x)/c(x)\ge1$.
We now argue for each algorithm.

\smallskip
\textit{PBGI.}
For each $x$ we defined a threshold $\alpha_{t}^{\mathrm{PBGI}}(x)$ by
$\mathrm{EI}_{f \mid x_{1:t}, y_{1:t}}\left(x; \alpha_{t}^{\mathrm{PBGI}}(x)\right)=c(x)$.
Since $\mathrm{EI}_{f \mid x_{1:t}, y_{1:t}}$ is increasing in its second argument,
\begin{equation}
\;y_{1:t}^*\;\ge\;\alpha_{t}^{\mathrm{PBGI}}(x)\;\Longleftrightarrow\;
\alpha_{t}^{\mathrm{EI}}(x)/c(x)\ge1.
\end{equation}
The existence of a point with ratio at least~$1$ therefore implies that the set
$
\mathcal S_{t}=\{x: y_{1:t}^*\ge \alpha_{t}^{\mathrm{PBGI}}(x)\}
$
is non-empty. PBGI chooses $x_{t+1}$ with the \emph{smallest} threshold,
and thus
\begin{equation}
\alpha_{t}^{\mathrm{EI}}(x_{t+1})=\mathrm{EI}_{f \mid x_{1:t}, y_{1:t}}\left(x_{t+1}; y_{1:t}^*\right)\ge\mathrm{EI}_{f \mid x_{1:t}, y_{1:t}}\left(x_{t+1}; \alpha_{t}^{\mathrm{PBGI}}(x_{t+1})\right)=c(x_{t+1}).
\end{equation}

\smallskip
\textit{(Log)EIPC.}
By definition $x_{t+1}$ maximizes $\log(\alpha_{t}^{\mathrm{EI}}(x)/c(x))$ and  $\alpha_{t}^{\mathrm{EI}}(x)/c(x)$, hence 
$\alpha_{t}^{\mathrm{EI}}(x_{t+1})\ge c(x_{t+1})$.

\smallskip
Thus for \emph{both} algorithms, we have
\begin{equation}
\label{eq:key-LB}
\alpha_{t}^{\mathrm{EI}}(x_{t+1})\ge c(x_{t+1}),
\qquad\text{for all }t < \tau.
\end{equation}
\end{proof}
We can now use \cref{lem:EI_lower_bound} to prove the following theorem, where we show that our PBGI/LogEIPC stopping rule (paired with the PBGI or LogEIPC acquisition function) also achieves cost-adjusted simple regret no worse than a naive baseline---stopping-immediately (\emph{Immediate}).

\ThmPositiveUtil*

\begin{proof}
We treat the two algorithms---meaning, the two acquisition function and stopping rule pairs---together and write
$
\mathcal F_{t}=\sigma\!\left(\{x_i,y_i\}_{i=1}^{t}\right)
$
for the filtration generated by the observations.
Since we are minimizing, the one–step improvement after
iteration~$t$ is
$y_{1:t-1}^*-y_{1:t}^*$, where recall that $y_{1:t}^* = \min_{1 \leq i \leq t} y_i$.
By our assumption about the random function $f$, 
$\E[y_{1:1}^*] = \mu(x_1)$ and
the quantity $y_{1:1}^*-\min\nolimits_{x\in X}f(x)$ has finite expectation $U<\infty$. Denote the posterior expected improvement function as $\alpha_{t}^{\mathrm{EI}}(x) = \mathrm{EI}_{f \mid x_{1:t}, y_{1:t}}\left(x; y^*_{1:t}\right)$. 


By \Cref{lem:EI_lower_bound}, $c(x_t)\le \alpha_{t-1}^{\mathrm{EI}}(x_t)$ for all $t\leq\tau$. Set
$
\Delta_t = y_{1:t-1}^* - y_{1:t}^* \ge 0
$
for $2\leq t\leq\tau$.
Conditional on $\mathcal F_{t-1}$ and the choice of $x_t$, we have
\begin{equation}
\E[\Delta_t \mid \mathcal F_{t-1},x_t]
      = \alpha_{t-1}^{\mathrm{EI}}(x_t).
\end{equation}      
Taking expectations and summing up from 2 to $\tau$, 
\begin{align}
\E\!\left[\sum_{t=2}^{\tau}\alpha^{\mathrm{EI}}_{t-1}(x_t)\right]
&= \E\!\left[\sum_{t=2}^{\infty}\mathbf{1}_{\{t\le \tau\}}\,
      \alpha^{\mathrm{EI}}_{t-1}(x_t)\right] \nonumber\\
&= \E\!\left[\sum_{t=2}^{\infty}\mathbf{1}_{\{t\le \tau\}}\,
      \E\!\left[\Delta_t \mid \mathcal F_{t-1}, x_t\right]\right] \nonumber
\\
&= \E\!\left[\sum_{t=2}^{\infty}\mathbf{1}_{\{t\le \tau\}}\,\Delta_t\right]       \quad(\text{tower property}) 
 \nonumber\\
&= \E\!\left[\sum_{t=2}^{\tau}\Delta_t\right]= \E\!\left[y^*_{1:1} - y^*_{1:\tau}\right]. \label{eq:max_total_improvment}
\end{align}

% expand on this part 
Summing \eqref{eq:key-LB} over $t\leq\tau$, taking expectations, and applying \eqref{eq:max_total_improvment}, we have
\begin{equation}
\mathbb{E}\left[\sum_{t=2}^{\tau}c(x_t)\right]
      \le
\mathbb{E}\left[\sum_{t=2}^{\tau}\alpha_{t-1}^{\mathrm{EI}}(x_t)\right]
      \le \E\!\left[y^*_{1:1} - y^*_{1:\tau}\right]. \label{eq:bound_on_cost} %\mu-\mathbb{E}\left[\min_{x\in X}f(x)\right].
\end{equation}
Adding the term $y^*_{1:\tau}$ on both sides of \eqref{eq:bound_on_cost} gives
\[
\mathbb{E}\!\left[y^*_{1:\tau}+\sum_{t=2}^{\tau} c(x_t)\right]
\le \mathbb{E}[\,y^*_{1:1}\,].
\]
Finally, adding $c(x_1)$ to both sides and subtracting $\mathbb{E}\!\left[\min_{x\in X} f(x)\right]$ yields
\begin{align*}
\mathbb{E}\!\left[y^*_{1:\tau}-\min_{x\in X} f(x)+\sum_{t=1}^{\tau} c(x_t)\right]
&\le
\mathbb{E}\!\left[y^*_{1:1}-\min_{x\in X} f(x)+c(x_1)\right]
\\
&=
\mathbb{E}\!\left[y_1-\min_{x\in X} f(x)+c(x_1)\right]
\\
&=
\mu(x_1) - \mathbb{E}\!\left[\min_{x\in X} f(x)\right] + \mathbb{E}[c(x_1)] \\
&= U+C,
\end{align*}
since $\mathbb{E}[y_1]=\mu(x_1)$, $C=c(x_1)$, and $U=\mu(x_1) - \mathbb{E}\left[\min_{x\in X} f(x)\right]$. This is exactly~\eqref{eqn:cost_adjusted_regret_bound}.
\end{proof}

\begin{remark}
    The quantities $U$ and $C$ depend on the choice of the initial point $x_1$. Choosing $x_1$ to minimize $\mu(x_1)+c(x_1)$ yields a tighter bound.
\end{remark}

\begin{remark}
In practice, one may center the objective by replacing $f(x)$ with $f(x) - \mu(x)$ 
and then using a zero-mean GP prior.  
This changes only the parameterization of the posterior mean (which can be shifted back by $\mu(x)$) 
and does not affect the implementation of the acquisition functions or stopping rules.
\end{remark}

Notably, this guarantee may not hold for other acquisition–stopping rule pairings. Moreover, in the worst case, this is the best guarantee we can hope for---for instance, the evaluation costs can be uniformly high and no point is worth evaluating.
We now formalize this worst-case regime.

\begin{proposition}[Immediate stopping is optimal under large costs]
\label{prop:immediate_optimal}
Suppose the objective function is bounded as $f(x) \in [a,b]$ for all $x \in \mathcal X$, and define $B := b-a$. Consider any Bayesian optimization policy that starts from an initial evaluation $x_1$ and may take additional evaluations. If the evaluation cost satisfies
\[
c(x) \ge B \quad \text{for all } x \in \mathcal X,
\]
then the policy that stops immediately after the initial evaluation is optimal with respect to expected cost-adjusted regret.
\end{proposition}

\begin{proof}
Let $\pi_{\mathrm{imm}}$ denote the policy that stops immediately after the initial evaluation $x_1$, and let $\pi$ be any other policy with stopping time $\tau \ge 1$. Recall that the cost-adjusted simple regret is
\[
R_c(\pi)
:=
\mathbb E\!\left[
\min_{1 \le t \le \tau} f(x_t) - \min_{x \in \mathcal X} f(x)
+ \sum_{t=1}^\tau c(x_t)
\right].
\]

Since $y_1 = f(x_1)$ and $f(x) \in [a,b]$ for all $x \in \mathcal X$, we have
\[
0 \le y_1 - \min_{x \in \mathcal X} f(x) \le b-a = B.
\]

The immediate-stopping policy has $\tau=1$, so
\[
R_c(\pi_{\mathrm{imm}})
=
\mathbb E\!\left[
y_1 - \min_{x \in \mathcal X} f(x)
+ c(x_1)
\right].
\]

Now consider an arbitrary policy $\pi$.

If $\tau = 1$, then $\pi$ coincides with immediate stopping and incurs the same cost-adjusted simple regret.

If $\tau \ge 2$, then the simple regret term is nonnegative, and the policy incurs at least one additional evaluation beyond $x_1$. Hence,
\[
\min_{1 \le t \le \tau} f(x_t) - \min_{x \in \mathcal X} f(x)
+ \sum_{t=1}^\tau c(x_t)
\ge
c(x_1) + c(x_2).
\]
Using $c(x) \ge B$ for all $x$,
\[
c(x_1) + c(x_2)
\ge
c(x_1) + B
\ge
c(x_1) + \bigl(y_1 - \min_{x \in \mathcal X} f(x)\bigr).
\]

Therefore, in all cases,
\[
\min_{1 \le t \le \tau} f(x_t) - \min_{x \in \mathcal X} f(x)
+ \sum_{t=1}^\tau c(x_t)
\ge
y_1 - \min_{x \in \mathcal X} f(x) + c(x_1).
\]
Taking expectations yields
\[
R_c(\pi) \ge R_c(\pi_{\mathrm{imm}}),
\]
which proves optimality of immediate stopping.
\end{proof}

\ThmTermination*
\begin{proof}
    Immediately by \cref{thm:non_neg_util} and the fact that $y^*_{1:\tau}-\min_{x\in X} f(x) \geq 0$ pointwise, we have that the expected cumulative cost of the algorithm 
    \begin{equation}
        \mathbb{E}\left[\sum_{t=1}^{\tau} c(x_t)\right] \leq U + C. 
    \end{equation}
    By Markov's inequality, for any $\delta \in (0, 1)$,
    \begin{align*}
        \Pr\left[\sum_{t=1}^{\tau} c(x_t) \leq \frac{U + C}{\delta} \right] \geq 1 - \delta. 
    \end{align*}
    Since $c(x_t)\geq c_0$ for all $t<\tau$, it follows that 
    \begin{align*}
        \tau = \sum_{t=1}^\tau 1 \leq  \sum_{t=1}^\tau \frac{c(x_t)}{c_0} = \frac{1}{c_0} \sum_{t=1}^\tau c(x_t).  
    \end{align*}
    Therefore, 
    \begin{align*}
        \Pr\left[\tau \leq \frac{U+C}{\delta \cdot c_0}\right] \geq   \Pr\left[\frac{1}{c_0} \sum_{t=1}^{\tau} c(x_t) \leq \frac{U + C}{\delta \cdot c_0} \right] = \Pr\left[\sum_{t=1}^{\tau} c(x_t) \leq \frac{U + C}{\delta} \right]  \geq 1 - \delta. 
    \end{align*}
\end{proof}

  
\ThmCosts*

\begin{proof}
Since the Bayesian optimization algorithm considers the post-scaling cost, by \cref{thm:non_neg_util} and the fact that $y^*_{1:\tau}-\min_{x\in X} f(x) > 0$ pointwise, we have 
\begin{equation}
\mathbb{E}\left[\sum_{t=1}^{\tau}\lambda \cdot c(x_t)\right] \leq \mathbb{E}\!\left[y^*_{1:\tau}-\min_{x\in X} f(x)+\sum_{t=1}^{\tau} \lambda \cdot c(x_t)\right] \le \mu(x_1)-\mathbb{E}\left[\min_{x\in X}f(x)\right] + \lambda C = U + \lambda C. 
\end{equation}

Since $\lambda = U/(B-C)$, we have \[\E\left[\sum_{t=1}^{\tau} c(x_t)\right] \leq \frac{U + \lambda C}{\lambda} = \frac{U}{\lambda} + C = B.\] 
\end{proof}

\begin{remark}
In the Bayesian regret setting where $f$ is drawn from a Gaussian process, i.e., $f \sim \mathcal{GP}(\mu(\cdot), K)$, the term $U = \mu(x_1) - \mathbb{E}\left[\min_{x\in X} f(x)\right]$ can be further bounded above using classical results on the expected supremum/infimum of Gaussian processes; see, for example, \citet[Theorem 10.1]{lifshits2012lectures}.
In many practical settings, the objective is naturally bounded. For example, in our AutoML experiments the objective is accuracy, which lies in $[0,100]$. This implies
\[
U = \mu(x_1) - \mathbb{E}\left[\min_{x\in X} f(x)\right] \le 100,
\]
providing a simple and interpretable upper bound on $U$.
\end{remark}
 

